% ============================================================================
% Fichier  : partie4/P04_C12_violations.tex
% Titre    : Les Violations de Données
% Auteur   : MINKA MI NGUIDJOÏ Thierry Emmanuel
% Version  : 1.0 -- novembre 2026
% Corrigés : corriges/Corr_C12.tex
% ============================================================================
% Sources : loi n° 2024/017, art. 5 (violation), 22, 63, 73 (texte
% officiel) ; ISO/IEC 27035 ; NIST SP 800-61 r2 ; Minka Mi Nguidjoï,
% Massoussi Songo et Bialo (2026, manuscrit) pour la tension T2 et les
% articles 102 et 154 du Code de procédure pénale cités par eux.
% Faits nouveaux (bible) : lettre de l'OPJ du 20/03/2025 ; aucune
% information des clients ; MBOLO, consultations de mai 2026.
% ============================================================================

\chapter{Les Violations de Données}
\label{chap:violations}
\chaptermark{Les violations de données}

\epigraph{%
    « Une violation tenue secrète protège l'organisation un temps,\\
    et laisse la personne sans défense. »%
}{Principe du cours}

\niveauChapitre{Intermédiaire}{%
    Chapitres~\ref{chap:securite} et~\ref{chap:loi}~; réponse à incident
    (UE MSN-2S3-SEC07)
}

\begin{boxobjectifs}
\begin{itemize}
    \item Reconnaître une violation de données au sens de l'article~5,
          y compris l'accès interne non autorisé.
    \item Appliquer l'obligation d'information « sans délai » de
          l'article~22.
    \item Qualifier une violation du point de vue des personnes.
    \item Auditer le dispositif de gestion des violations et le registre
          qui les documente.
    \item Analyser la rencontre entre l'obligation d'informer et le secret
          d'une enquête pénale.
\end{itemize}
\end{boxobjectifs}

\bigskip

% ============================================================================
\section*{Avant de commencer : une lettre de l'OPJ}
% ============================================================================

Bafoussam, vendredi 21~août~2026. Dans le dossier de l'incident de
mars~2025, Mme~BELLO trouve une lettre de l'adjudant-chef NANA, officier
de police judiciaire, datée du 20~mars~2025~: « Dans l'intérêt de
l'enquête, il vous est demandé de ne pas informer vos clients de la
compromission avant instruction contraire. » Aucune instruction contraire
n'a été donnée. Le jugement a été rendu en février~2026. Le 21~août, les
18\,412 personnes dont les dossiers ont été copiés n'ont toujours reçu
aucune information.

Le directeur général s'en tient à la lettre~: « La police nous a demandé de
nous taire. » L'auditrice note que la lettre visait l'enquête, que
l'enquête est close, et que personne n'a vérifié depuis si la demande
tenait toujours.

\begin{boxalert}
\textbf{Violation typique.} Traiter une violation comme un problème
technique à corriger en silence, ou s'abriter indéfiniment derrière une
demande de discrétion devenue sans objet.

\smallskip
\textbf{Règle~:} une violation s'informe « sans délai » à l'Autorité et aux
personnes (article~22). Une raison de différer, si elle existe, se
documente, se limite dans le temps et se réexamine~; quand elle cesse,
l'information est due.
\end{boxalert}

% ============================================================================
\section{Qu'est-ce qu'une Violation}
\label{sec:c12-definition}
% ============================================================================

\begin{boxdef}
\termecle{Violation de données à caractère personnel}~: selon l'article~5,
« toute atteinte à la sécurité des données à caractère personnel
entraînant, de manière accidentelle ou illicite, la destruction, la perte,
l'altération, la divulgation non autorisée desdites données transmises,
conservées ou traitées d'une autre manière, ou l'accès non autorisé à de
telles données ».
\end{boxdef}

La définition est large. Une violation n'est pas forcément une attaque~:
une erreur d'envoi, un ordinateur portable perdu, une sauvegarde effacée
par erreur en sont. Elle n'est pas forcément externe~: l'accès non autorisé
d'un agent à un dossier, comme à MBOLO, en est une. Elle n'est pas
forcément une fuite~: la destruction ou l'indisponibilité de données
indispensables en est une aussi.

\begin{tableaucours}{L{3.0cm}L{4.6cm}L{4.4cm}}{%
    \tableauHeader{Atteinte} &
    \tableauHeader{Exemple} &
    \tableauHeader{Risque pour la personne}
}
Confidentialité & Fichier clients copié, dossier consulté sans raison,
courriel au mauvais destinataire & Fraude, chantage, discrimination,
atteinte à la réputation \\
Intégrité & Résultat d'analyse attribué au mauvais patient & Soin inadapté,
décision erronée \\
Disponibilité & Dossiers chiffrés par un rançongiciel, sauvegardes
détruites & Soin retardé, droits impossibles à exercer \\
\end{tableaucours}
\vspace{-0.8em}
\begin{center}{\small\itshape Tableau~12.1. Les trois atteintes et leurs effets}\end{center}

% ============================================================================
\section{L'Obligation d'Informer}
\label{sec:c12-article22}
% ============================================================================

\begin{boxjuris}
\textbf{\loicam{n°~2024/017}, article~22}~\cite{loi2024017}~:
\textit{« (1) Dès connaissance de toute violation des données à caractère
personnel, le responsable du traitement ou le sous-traitant est tenu
d'informer sans délai, l'Autorité de protection des données à caractère
personnel et la personne concernée. (2) En tout état de cause, et sans
préjudice des dispositions de l'alinéa 1 ci-dessus, le responsable du
traitement ou le sous-traitant est tenu de mettre en œuvre les mesures de
protection techniques et organisationnelles appropriées, telles que le
chiffrement, l'instauration des autorisations restrictives d'accès aux
données. »}

\smallskip
\textbf{Portée.} L'obligation est plus large que dans le droit européen sur
deux points~: elle vise \textit{toute} violation, sans seuil de risque, et
elle impose d'informer la personne dans tous les cas. Elle est moins
précise sur un troisième~: « sans délai » n'est pas chiffré, et le contenu
de l'information n'est pas fixé.
\end{boxjuris}

En l'absence de contenu fixé, le cours retient une information minimale,
déclarée comme critère~: la nature de la violation, les catégories et le
nombre approximatif de personnes et de données touchées, les conséquences
probables pour les personnes, les mesures prises et celles que la personne
peut prendre elle-même (changer un mot de passe, surveiller ses comptes),
et un contact. Pour « sans délai », l'audit mesure le temps écoulé entre la
connaissance de la violation et chaque information, et le rapproche des
raisons documentées d'un éventuel retard.

\begin{boxinfo}
\textbf{Le sous-traitant aussi.} L'article~22 vise « le responsable du
traitement ou le sous-traitant ». Le contrat doit organiser qui informe
qui, et dans quel ordre~: un prestataire qui découvre une violation doit
au moins prévenir immédiatement son client (clause~9 du
Chapitre~\ref{chap:soustraitants}).
\end{boxinfo}

% ============================================================================
\section{Gérer une Violation}
\label{sec:c12-gerer}
% ============================================================================

\begin{boxprocedure}
\textbf{Auditer le dispositif de gestion des violations.}
\begin{enumerate}
    \item Vérifier qu'une procédure écrite existe, connue des équipes et des
          sous-traitants, avec un point de contact unique.
    \item Vérifier la détection~: comment l'organisation apprend-elle une
          violation (alertes, revues de journaux, signalements)~?
    \item Obtenir le registre des violations~; à défaut, le reconstituer à
          partir des incidents de sécurité, des plaintes et des sinistres.
    \item Pour chaque violation, relever les dates de connaissance, de
          confinement, d'information de l'Autorité et des personnes.
    \item Vérifier la qualification~: données, personnes, conséquences
          pour elles.
    \item Vérifier le contenu des informations envoyées.
    \item Vérifier les enseignements tirés et les mesures correctives.
\end{enumerate}
\end{boxprocedure}

La loi n'impose pas explicitement un registre des violations. Il découle
cependant de deux obligations~: prouver qu'on a informé « sans délai »
suppose de dater la connaissance~; et la traçabilité de l'article~27
suppose de pouvoir reconstituer ce qui s'est passé. Un registre où chaque
violation, même non notifiée, est décrite avec sa date, sa qualification
et sa décision, est la seule manière de le démontrer.

% ============================================================================
\section{Informer pendant une Enquête Pénale}
\label{sec:c12-penal}
% ============================================================================

L'affaire de NYANGA pose une question que la loi ne règle pas. L'article~22
impose d'informer sans délai les personnes~; la procédure pénale protège
le secret de l'enquête. Dans leur travail en cours, l'auteur de ce cours et
ses coauteurs relèvent que le Code de procédure pénale impose le secret de
la procédure durant l'enquête de police judiciaire (article~102) et le
secret de l'information judiciaire (article~154)~\cite{minka2026chainepenale}.
Une information des personnes peut révéler l'existence de l'enquête,
l'identité d'un suspect ou une opération en cours.

Les auteurs qualifient cette rencontre de « tension procédurale », et non
d'incompatibilité~: les deux obligations peuvent être conciliées si l'on
organise le moment et les modalités de l'information. Ils proposent qu'un
mécanisme permette au magistrat compétent d'apprécier ce moment lorsque
l'information risque de compromettre l'enquête. En attendant un tel
mécanisme, l'auditeur vérifie quatre choses~:

\begin{enumerate}
    \item la demande de différer existe par écrit, émane d'une autorité
          compétente et vise l'enquête~;
    \item l'Autorité de protection a été informée, ou l'organisation a
          documenté pourquoi elle ne pouvait pas l'être~;
    \item l'organisation a réexaminé la demande à intervalles réguliers~;
    \item l'information a été donnée dès que la raison de différer a cessé.
\end{enumerate}

S'y ajoute une question de temps, que les mêmes auteurs nomment « tension
temporelle »~: la violation de NYANGA date du 14~mars~2025, après la
promulgation de la loi mais pendant le délai de mise en conformité de
l'article~73. L'obligation d'informer s'appliquait-elle déjà~? L'auditeur
ne tranche pas cette question de droit~; il constate en revanche qu'au
24~juin~2026, à l'échéance du délai, la violation n'était pas réparée du
point de vue des personnes, et qu'aucune raison documentée ne justifiait
plus le silence. Le Chapitre~\ref{chap:recherche} présente ces deux
tensions comme des pistes de recherche.

% ============================================================================
\section{Les Violations de NYANGA et de MBOLO}
\label{sec:c12-cas}
% ============================================================================

\begin{boxscenario}{Missions NYANGA et MBOLO. Deux violations, deux traitements}
\textbf{NYANGA.} Connaissance~: 14~mars~2025, 9~heures. Qualification
technique faite par l'enquête (3,2~Go, 18\,412 fichiers). Aucune
qualification du point de vue des personnes. Demande de l'OPJ du
20~mars~2025, jamais réexaminée. Jugement le 15~février~2026. Au
31~août~2026, aucune information des personnes ni de l'Autorité~; aucun
registre des violations.

\textbf{MBOLO.} Connaissance~: 13~mai~2026, par la plainte de la patiente.
Investigation le jour même (journal des consultations), comptes désactivés,
procédure disciplinaire. Patiente informée par écrit le 15~mai, avec la
liste des consultations non justifiées et les mesures prises. L'Autorité
n'a pas été informée~; la directrice a noté au registre, créé le 14~mai,
qu'aucun interlocuteur n'était identifié.

\textbf{Analyse} (confiance élevée). MBOLO a traité sa violation comme la
loi le demande pour ce qui dépendait de lui, et l'a documentée. NYANGA a
traité la sienne comme un dossier pénal, sans jamais la regarder comme une
violation de données envers ses clients.
\end{boxscenario}

\begin{boxlecture}{Le domaine D5, violations, à NYANGA}
\begin{tableaucours}{L{7.8cm}cL{3.2cm}}{%
    \tableauHeader{Exigence} &
    \tableauHeader{Cote} &
    \tableauHeader{Preuve}
}
Une procédure de gestion des violations existe & \NON{} & Aucune \\
Un registre des violations est tenu & \NON{} & Aucun \\
La violation de 2025 est qualifiée pour les personnes & \NON{} & Rapport
technique seul \\
La demande de différer est écrite et motivée & \OUI{} & Lettre du
20~mars~2025 \\
La demande a été réexaminée & \NON{} & Aucune démarche \\
Les personnes ont été informées quand la raison a cessé & \NON{} & Rien au
31~août~2026 \\
L'Autorité a été informée & \NEK{} & Autorité non organisée, aucune
démarche documentée \\
L'obligation d'informer s'appliquait en mars~2025 & \NEK{} & Question de
droit non tranchée \\
\end{tableaucours}

\medskip
\textbf{Décompte.} Un \OUI{}, cinq \NON{}, deux \NEK{} sur huit.
\hfill\lecture{0{,}125}{0{,}25}{1{,}30}

\medskip
\textbf{Lecture.} Les deux \NEK{} ne changent pas la recommandation~: quelle
que soit la réponse à la question de droit, la raison de se taire a cessé
en février~2026, et l'information des 18\,412 personnes est due. C'est la
première recommandation du rapport.
\end{boxlecture}

% ============================================================================
\section{Ce que les Violations demandent à chaque Métier}
\label{sec:c12-metiers}
% ============================================================================

\begin{boxmetier}{Responsable de traitement}
Désigner un point de contact pour les violations, décider de l'information
et la signer. Réexaminer par écrit toute demande de discrétion et lever le
silence dès qu'elle cesse.
\end{boxmetier}

\begin{boxmetier}{RSSI}
Alimenter le registre des violations depuis les incidents de sécurité, avec
les éléments techniques, et apporter la qualification du point de vue des
données touchées.
\end{boxmetier}

\begin{boxmetier}{Juriste conformité}
Articuler l'obligation d'informer avec les demandes des autorités
judiciaires~: obtenir des demandes écrites, limitées dans le temps, et
documenter chaque décision.
\end{boxmetier}

\begin{boxresume}
\begin{itemize}
    \item Une violation est toute atteinte à la confidentialité, à
          l'intégrité ou à la disponibilité de données personnelles,
          accidentelle ou illicite, interne ou externe.
    \item L'article~22 impose d'informer sans délai l'Autorité et la
          personne, pour toute violation.
    \item Un registre des violations est la seule manière de prouver le
          respect du « sans délai ».
    \item L'information peut rencontrer le secret de l'enquête pénale~: une
          demande de différer se documente, se limite et se réexamine.
    \item À NYANGA, 18\,412 personnes n'avaient toujours pas été informées
          en août~2026~; MBOLO a informé sa patiente en deux jours.
\end{itemize}
\end{boxresume}

\begin{boxinfo}
\textbf{Points de contrôle.} Vous avez compris ce chapitre si vous savez~:
\begin{itemize}
    \item reconnaître une violation interne~;
    \item rédiger l'information d'une personne après une violation~;
    \item auditer une demande de différer l'information.
\end{itemize}
\end{boxinfo}

\begin{boxexo}[title={\ding{45}\quad Exercice~12.1. Violation ou non \niveaubase}]
(a) un relevé de compte envoyé au mauvais client~; (b) une tentative
d'intrusion bloquée par le pare-feu~; (c) un serveur indisponible deux
jours, sans perte de données~; (d) un portable chiffré perdu dans un
taxi~; (e) un portable non chiffré volé.
\end{boxexo}

\begin{boxexo}[title={\ding{45}\quad Exercice~12.2. Informer M. TAGNE \niveaumoyen}]
Rédigez l'information que NYANGA doit adresser aux 18\,412 personnes,
selon le contenu minimal du cours. Que dites-vous du retard~?
\end{boxexo}

\begin{boxexo}[title={\ding{45}\quad Exercice~12.3. La chronologie \niveaumoyen}]
À partir des dates de la section~\ref{sec:c12-cas}, calculez pour NYANGA le
délai entre la connaissance de la violation et le jugement, puis entre le
jugement et le 31~août~2026. Rédigez la phrase du rapport sur le respect du
« sans délai ».
\end{boxexo}

\begin{boxexo}[title={\ding{45}\quad Exercice~12.4. La demande de l'OPJ \niveauexpert}]
Un OPJ demande aujourd'hui à MBOLO de ne pas informer des patientes dont les
dossiers ont été consultés par un agent sous enquête. Rédigez la réponse de
MBOLO à l'OPJ et la note au registre des violations, en organisant la
conciliation entre l'article~22 et le secret de l'enquête.
\end{boxexo}

\bigskip

\begin{boxinfo}
\textbf{Pour aller plus loin.}
\begin{itemize}
    \item L'ISO/IEC~27035~\cite{iso27035} et le NIST SP~800-61~\cite{nist80061r2},
          pour la gestion des incidents.
    \item Le cours de forensique de l'auteur (UE MSN-2S3-SEC07), pour
          l'enquête sur l'incident de NYANGA.
\end{itemize}
\end{boxinfo}

\begin{boxinfo}
\textbf{Transition.} Les constats s'accumulent~: registre, fondements,
droits, sécurité, sous-traitants, transferts, violations. Il reste à les
transformer en rapport et en programme de mise en conformité. C'est
l'objet du Chapitre~\ref{chap:programme}.
\end{boxinfo}
